\originalsection{条件记忆怎样分开容量与计算}\label{sec:engram-bridge}
我们已经见过 embedding 的按行读取：词表很大，每个 token 只读取一行。Engram\cite{engram}把这个思路推广到压缩后的局部 $N$-gram，用多头确定性哈希查找可训练向量，再依据当前隐藏状态门控和融合。与 MoE 的条件计算相比，它探索条件记忆这个额外分配轴；与 RAG 相比，这里读取的是模型内学得的表项，不是带文档出处的外部检索结果。两者都叫“记忆”时，应先问读出的对象与更新过程。

下面是解释机制的简化数学模型，不替代原论文的多分支和卷积公式。令 $u_t$ 为离散 token，$H$ 为哈希函数，$z_t=H(u_{t-n+1:t})$，取表向量 $e_t=E[z_t]\in\R^m$。隐藏状态记为 $a_t\in\R^d$，$V\in\R^{d\times m}$，门控 $g:\R^d\times\R^m\to\R^d$，并设 $a'_t=a_t+g(a_t,e_t)\odot Ve_t$。查找地址依赖离散 token，可在处理序列前计算；表向量、投影和连续门控仍通过训练得到。若一个 batch 多次读同一行，梯度必须累加：
\[
\frac{\partial L}{\partial E[j]}=\sum_{t:z_t=j}\frac{\partial L}{\partial e_t}.
\]
这与数学工具附录的共享参数完全相同。哈希碰撞使不同上下文共享表项；多头可缓解碰撞，但不能保证语义唯一。门控使同一静态表项在不同动态上下文中产生不同影响，又不把查找变成可靠事实数据库。

\texttt{research\_bridges.py} 对 17 行三维表做真实有限差分；5 个样本只访问 2、7、11 行，其中 2 号访问三次。最大梯度误差为 $4.95\times10^{-11}$，未访问行梯度为零。另一个确定整数键例子将 1,000 个键单头映射到 257 个地址，双头联合签名有 1,000 种。后者不是一般无碰撞保证：模表大小、键域和哈希函数都由例子指定，两个表的总存储也增加了。

原论文比较固定参数与 FLOPs 的配置，并做分配比例和组件消融。读这些结果时，须分别看总参数、活跃计算、GPU 与主存、通信和吞吐。$O(1)$ 指某种地址查找随表项数的渐近成本，不表示整段 LLM 推理为常数时间。表容量增加还会消耗存储和带宽；作者特定硬件的主存卸载结果不等于普通 CPU 无代价复现。CKA 和 LogitLens 的层间相似性支持一种机制假说，但不足以识别唯一因果路径。配套全译保留原文摘要与表格的个别数字差异。

\originalsection{推理预算怎样转成下一轮监督}\label{sec:cat-bridge}
测试时计算花在当前问题上；训练则改变以后问题的回答分布。CaT\cite{cat}连接这两件事：从同一策略生成多条轨迹，聚合为伪参考，再将与伪参考的符合程度转成强化学习奖励。在难以写程序验证的任务中，它还用生成的细分规则和 LLM 评分器给奖励。这增加了可生成反馈的范围，也引入了聚合器、规则和评分器的偏差。

将全过程写为 $o_{1:G}\sim\pi_\theta$，$s=A(q,o_{1:G})$，$r_i=R(q,o_i,s)$，然后更新 $\theta$。$s$ 不是隐藏的标准答案；它也是一个推断结果。若所有轨迹共享误解，增加 $G$ 未必增加独立证据。若使用另一模型聚合，知识可能来自那个模型；“没有人工参考标注”不能简化成系统从无知识状态自行产生真理。

\begin{proposition}
二元答案的每条独立轨迹出错概率为 $p<1/2$，奇数 $G$ 的多数投票错误概率为
\[
\sum_{j=(G+1)/2}^{G}{G\choose j}p^j(1-p)^{G-j}.
\]
若同一组全部轨迹共享一个错误 Bernoulli 变量，则多数投票错误概率仍为 $p$，与 $G$ 无关。
\end{proposition}
\begin{proof}
独立时错误数为二项分布，多数错误对应其右尾。完全共享时，每组要么全部正确，要么全部错误，投票不改变事件。
\end{proof}
我们的真实 CPU 模拟设 $p=0.3$、每种条件 20,000 组。$G=9$ 时，独立错误的伪参考错误率为 $0.10055$，共享错误为 $0.30045$；$G=33$ 时分别为 $0.00795$ 与 $0.29570$。共享条件总是全体一致，仍经常错误。因此输出收敛能减少训练中的差异信号，却不意味着知识正确。这是人为规定相关性的机制实验，没有 LLM、医疗任务或 CaT 的真实评分器，不用于估计论文医疗可靠性。

作者在指定模型与 HealthBench、MATH-500 上报告效果及推理成本收益。学习者应查看：评价是否独立于训练规则、聚合与更新分别贡献多少、不同 judge 是否改变排序、预算节省是训练后单次推理还是包含训练总成本。一个有用的本地后续研究是在已知真值的任务上故意注入共享偏差，用外部留出规则审计伪参考；在开放域没有真值时，要把人工/工具证据的新增信息来源明确写出。

\originalsection{形式核验把研究闭环推进到哪里}\label{sec:alphaproof-bridge}
AlphaProof\cite{alphaproof}将形式数学状态、策略与价值学习、证明搜索以及测试时强化学习结合。它与 FunSearch 的共同点是把模型提议放在可检查的闭环中；区别是核验的对象为形式证明而不是只在若干输入上的程序分数。正式论文的奥赛结果还涉及人工形式化题目和与 AlphaGeometry 2 的联合贡献，应按系统和时间预算读，不能将联合成绩都记到一个模型。

设根目标为命题 $P$。一个动作把它变成子目标集合 $\{P_1,\ldots,P_k\}$；要接受该动作，需要每个子目标都有有效证明，这是 AND 条件。对同一目标尝试多个动作，只须其中一个成功，这是 OR 条件。若仅把子目标的价值均值当成已经证明，可能忽略一个尚未解决的分支；启发式价值用于排序，不是证书。搜索结束仍要由证明核检查整个构造。

\begin{example}
证明 $A\land B$ 的候选先给出了 $A$ 的证明，价值网络估计 $B$ 成功概率为 $0.99$。这不允许核验器接受 $A\land B$。若有两个候选，各自只完成一个不同分支，也不能仅凭平均价值宣布成功；需要组合成同一条合法证明并检查依赖和假设。
\end{example}
形式证书强于经验测试，但也有明确范围：核验的是形式命题在指定公理与环境中的成立；自然语言题意到形式命题的对应、模型的数据污染、证明搜索成本与新颖性仍需要外部审计。AI 辅助研究的分工因此是提出猜想与候选、搜索和核验、解释与独立审查，而不是把流畅叙述或高价值估计直接当证明。

\begin{exercise}
上面的条件查找例子，若相同表项在两次使用中收到相反梯度，其总梯度为零。能否说该表项从未被使用？
\end{exercise}
\begin{solution}
不能。梯度是两次依赖的和，可能抵消；它与没有计算路径不同。应检查具体访问和逐次作用，必要时设计输入干预。
\end{solution}
\begin{exercise}
多数投票随样本数增加更准确，需要哪些假设？若评分器总给多数答案奖励，为什么这个训练目标可能固化错误？
\end{exercise}
\begin{solution}
这里的证明需要二元、独立同分布且每样本错误概率小于一半；相关、系统性误解或不同分布都破坏推论。错误多数会被当成伪参考，奖励再提高其概率，从而减少正确少数的出现。必须有外部证据或独立评价来识别这种闭环偏差。
\end{solution}
